% =========================================================================
% Evaluation - TempoMem (restructured to follow the experiments layout of
% CST-ViT, Wang et al., IEEE IoT-J 2025: A. Settings and Details,
% B. Evaluating Details, C. Experimental Results, D. Comparison with SOTA,
% E. Ablation Experiment (incl. runtime/memory), F. Secondary task).
% Requires: booktabs. Tables use \resizebox{\columnwidth}{!}{...} where
% they would otherwise overflow the AAAI column (needs graphicx).
%
% Placeholders for numbers not yet measured are marked [TBD]. Search for
% "TBD" before submission.
% =========================================================================

\section{Experiments}
\label{sec:eval}

We evaluate TempoMem on two representative video understanding tasks:
video action recognition with ViViT-B on Kinetics-400, where the
prediction is read from a class token, and video object detection with
ViTDet-B on ImageNet VID, where the prediction depends on the full spatial
token grid. The two regimes stress the method in opposite ways, and we
report both the positive and the negative findings: content-based reuse is
highly effective for class-token readouts and is currently dominated by
per-position gating on dense spatial prediction. The comparison with the
state of the art analyzes why.

\subsection{Experimental Settings and Details}
\label{sec:setup}

\paragraph{Dataset and evaluation setup.}
For action recognition we use Kinetics-400 \cite{k400} and, for
backbone-scaling experiments, UCF101 \cite{ucf101}; both report Top-1 and
Top-5 accuracy. For object detection we use ImageNet VID \cite{imagenet},
reporting mAP@50. Compute is reported as GFLOPs, measured with in-code
operation counters (instrumented linear and matmul modules) rather than
analytic estimates, so all reductions reflect the operations actually
executed, including matching overhead. Latency is CUDA-event wall-clock
time per frame and memory is peak GPU allocation, both on a single NVIDIA
Quadro RTX 6000 (24\,GB) under a shared warm-up and reset protocol.

\paragraph{Model design for video action recognition.}
We use ViViT-B \cite{vivit} with factorized spatio-temporal attention,
16-frame clips at $224 \times 224$, tubelet $2 \times 16 \times 16$, and a
class-token readout; TempoMem replaces computation in the spatial module,
which dominates cost. For backbone scaling we additionally use VideoMAE
\cite{videomae} with ViT-S/B/L/H encoders.

\paragraph{Model design for video object detection.}
We use ViTDet-B \cite{vitdet} on ImageNet VID at $672 \times 672$. Its
windowed attention with relative-position biases means TempoMem runs in
the KV-reuse-only variant on windowed blocks, as described in the
Methodology.

\paragraph{Comparison methods and parameter ranges.}
\emph{Dense} is the unmodified backbone. \emph{Eventful} \cite{eventful}
recomputes the top-$k$ most-changed tokens per frame and reuses buffered
results for the rest; $k$ is its compute budget. \emph{STGT} \cite{stgt}
is a spatio-temporal token-gating baseline. \emph{MaskVD} \cite{maskvd}
skips background tokens using detection masks. \emph{ToMe} \cite{tome} is
a single-frame token-merging method that we also compose with TempoMem on
foreground tokens. We explore the accuracy--efficiency tradeoff by sweeping
the reuse fraction $\gamma \in \{0, 0.25, 0.5, 0.75\}$ (the fraction of
Eventful's recompute set that TempoMem drops), the prototype merge
schedule $T \in \{2, 4, 8\}$ passes, and the Eventful budget
$k \in \{128, \dots, 1024\}$. Here, $\gamma$ controls how aggressively
matched tokens are reused, while $T$ determines how coarsely the prototype
cache compresses history, jointly serving as the key parameters for
balancing accuracy and computational cost.

\subsection{Evaluating Details}
\label{sec:eval-details}

Each video runs the TempoMem pipeline: the first $W$ frames are caching
frames ($W = 4$ for ViTDet, the leading frames of each clip for ViViT),
and all later frames are matching frames. Unless a table says otherwise,
accuracy and GFLOPs are computed \emph{on matching frames only}, for every
method, with baselines evaluated on the same frames; this is the
steady-state accounting defined in the Methodology, and warm-up cost is
reported separately where relevant. Whole-stream numbers (all frames,
warm-up included) are labeled explicitly, and \textbf{we never compare a
matching-only number against a whole-stream number in the same table}.

It is important to note that TempoMem is a training-free, inference-time
optimization: the prototype cache and matching modules introduce no
learnable parameters and are only active during inference. All weights are
frozen, pretrained checkpoints, and every compared method is applied as a
training-free, inference-time modification, which isolates the effect of
the reuse mechanism itself.

\subsection{Experimental Results}
\label{sec:eval-results}

\paragraph{Standalone TempoMem.}
Table~\ref{tab:standalone-vivit} evaluates standalone TempoMem
(saliency-split background caching, KV reuse, and token substitution) on
ViViT-B/Kinetics-400. TempoMem removes 27\% of FLOPs with no Top-1 loss;
composing ToMe on the foreground pushes the reduction to 71\% at a
7.3-point Top-1 cost, showing that the two mechanisms stack but that
aggressive foreground reduction is where accuracy is spent. The Top-5 drop
under substitution (97 to 87) indicates that prototype substitution
perturbs the tail of the logit distribution more than its mode.

\begin{table}[t]
\centering
\small
\begin{tabular}{lccc}
\toprule
Method & GFLOPs & Top-1 (\%) & Top-5 (\%) \\
\midrule
ViViT-B (dense) & 537.6 & 71.0 & 97.0 \\
+ TempoMem & 392.0 ($-$27\%) & 71.0 & 87.0 \\
+ TempoMem + ToMe & 155.4 ($-$71\%) & 63.7 & 78.0 \\
\bottomrule
\end{tabular}
\caption{Standalone TempoMem on ViViT-B / Kinetics-400 (per-clip GFLOPs,
matching frames scored). TempoMem alone preserves Top-1 exactly at a 27\%
FLOPs reduction; adding ToMe on foreground tokens trades accuracy for a
much deeper reduction.}
\label{tab:standalone-vivit}
\end{table}

\paragraph{TempoMem on Eventful transformers.}
We next apply TempoMem as a recompute-set filter on top of Eventful
(Algorithm~\ref{alg:eventfultm}): Eventful's qkv gate proposes its top-$k$
changed tokens, and TempoMem drops the fraction $\gamma$ of them that best
match the prototype cache. With $\gamma = 0$ the pipeline is exactly
Eventful under our protocol, which gives a perfectly controlled comparison:
every row of Table~\ref{tab:eventful-vivit} differs from the first row
only in $\gamma$.

\begin{table}[t]
\centering
\small
\begin{tabular}{lccc}
\toprule
$\gamma$ & Top-1 (\%) & Top-5 (\%) & GFLOPs/frame \\
\midrule
0.00 (= Eventful) & 61.0 & 82.0 & 27.2 \\
0.25 & 59.0 & 79.0 & 20.6 \\
0.50 & 59.0 & 79.0 & 13.9 \\
0.75 & 59.0 & 79.0 & \textbf{7.2} \\
\bottomrule
\end{tabular}
\caption{Eventful + TempoMem on ViViT-B / Kinetics-400 (matching-frame
accounting for all rows). Accuracy is flat for all
$\gamma \ge 0.25$ while compute falls $3.8\times$: TempoMem can discard up
to 75\% of Eventful's recompute set for a fixed 2-point Top-1 cost.
Whole-clip reference points under the same setting: dense 73.0\% Top-1 at
3355 GFLOPs/clip; canonical Eventful 71.0\% at 613 GFLOPs/clip (not
comparable to the matching-only rows above).}
\label{tab:eventful-vivit}
\end{table}

\paragraph{Reading the accuracy--compute curve
(Fig.~\ref{fig:eventful-curve}).}
Figure~\ref{fig:eventful-curve} plots Top-1 against matching-frame
GFLOPs as $\gamma$ sweeps from 0 to 0.75. Two properties matter. First,
the curve is a step, not a slope: the entire accuracy cost (2 points) is
paid at the smallest nonzero $\gamma$, after which accuracy is exactly flat
at 59.0/79.0 while compute falls from 20.6 to 7.2 GFLOPs/frame. This says
the tokens TempoMem drops beyond the first quartile were genuinely
redundant: Eventful flagged them as changed, but the change was already
represented in the prototype cache. Second, the step localizes the damage:
the loss comes from the first tranche of dropped tokens (those nearest the
matching threshold), suggesting an adaptive $\gamma$ could recover it.
Prototype granularity is irrelevant here: merge schedules of $T \in
\{2,4,8\}$ passes produced identical Top-1, Top-5, and FLOPs at every
$\gamma$, so on classification the filter is robust to how finely the cache
is compressed.

\begin{figure}[t]
\centering
% [TBD] Insert plot: x = matching-frame GFLOPs/frame, y = Top-1 (%),
% one curve for Eventful+TempoMem (gamma sweep, gamma=0 endpoint is exact
% Eventful), horizontal reference line for dense Top-1 under the same
% matching-frame protocol. Source: results/comparison/plot_vivit_runs.log
\fbox{\parbox{0.9\columnwidth}{\centering\vspace{2em}
\textbf{[TBD: accuracy vs.\ compute curve]}\vspace{2em}}}
\caption{Top-1 accuracy vs.\ matching-frame GFLOPs for Eventful + TempoMem
on ViViT-B / Kinetics-400 as the reuse fraction $\gamma$ increases (left to
right: $\gamma = 0.75, 0.5, 0.25, 0$). Accuracy is flat for all
$\gamma \ge 0.25$; the full accuracy cost is paid at the first step.}
\label{fig:eventful-curve}
\end{figure}

\subsection{Comparison Experiments with SOTA Methods}
\label{sec:eval-vid}

This subsection verifies where TempoMem stands relative to Eventful
\cite{eventful}, STGT \cite{stgt}, and MaskVD \cite{maskvd} on ImageNet
VID, under two accountings kept strictly separate.

\paragraph{Controlled comparison.}
The apples-to-apples evidence is Table~\ref{tab:vitdet-sweep}: Eventful and
Eventful + TempoMem under the identical matching-frame protocol, differing
only in $\gamma$ (Eventful budget $k = 512$). Unlike classification,
detection accuracy falls monotonically with $\gamma$ (90.3\% to 59.3\%
mAP@50) as compute falls from 60.3 to 19.3 GFLOPs/frame; at $\gamma = 0.5$,
halving Eventful's compute costs 13 mAP points. Sweeping the Eventful
budget $k \in \{128, \dots, 1024\}$ shows the same picture along the whole
curve: at matched compute, Eventful reaches higher mAP@50 than Eventful +
TempoMem, and at matched mAP@50 Eventful is cheaper. \textbf{On detection,
the recompute-set filter does not sit on Eventful's Pareto front}; we
report this negative result rather than tuning around it, and analyze its
cause below.

\begin{table}[t]
\centering
\small
\begin{tabular}{lccc}
\toprule
$\gamma$ & merge passes $T$ & mAP@50 (\%) & GFLOPs/frame \\
\midrule
0.00 (= Eventful) & -- & 90.3 & 60.3 \\
0.25 & 2 & 86.9 & 46.6 \\
0.25 & 4 & 86.4 & 46.6 \\
0.25 & 8 & 85.2 & 46.6 \\
0.50 & 2 & 77.8 & 32.9 \\
0.50 & 4 & 76.9 & 32.9 \\
0.50 & 8 & 74.7 & 32.9 \\
0.75 & 2 & 59.3 & 19.3 \\
\bottomrule
\end{tabular}
\caption{Eventful + TempoMem on ViTDet-B / ImageNet VID
($k = 512$, matching-frame accounting for all rows). Detection accuracy
falls monotonically with the reuse fraction $\gamma$, and coarser
prototypes (more merge passes) hurt, the opposite of the classification
result.}
\label{tab:vitdet-sweep}
\end{table}

\paragraph{Reference frontier.}
Table~\ref{tab:vitdet-sota} places the compared methods under whole-stream
per-frame accounting (all frames, warm-up included), together with their
measured latency and peak memory. Eventful and STGT define the current
frontier; MaskVD follows. The Eventful + TempoMem row requires a
whole-stream re-measurement and is marked \textbf{[TBD]}; from the
matching-only evidence above we expect it below the Eventful frontier on
this task.

\begin{table}[t]
\centering
\resizebox{\columnwidth}{!}{%
\begin{tabular}{lcccc}
\toprule
Method & GFLOPs/fr. & Latency (ms) & Mem.\ (MB) & mAP@50 (\%) \\
\midrule
ViTDet-B (dense) & 174.5 & 69.6 & 760 & 84.1 \\
Eventful & 60.5 & 58.3 & 2118 & 84.0 \\
STGT & 68.5 & 56.2 & 1242 & 83.8 \\
MaskVD & 88.8 & 56.0 & 927 & 83.1 \\
TempoMem (standalone) & 156.5 & 101.8 & 929 & 72.0 \\
Eventful + TempoMem & \textbf{[TBD]} & \textbf{[TBD]} & \textbf{[TBD]} &
  \textbf{[TBD]} \\
\bottomrule
\end{tabular}}
\caption{ImageNet VID comparison with the state of the art under
whole-stream per-frame accounting and an identical GPU timing protocol.
All rows share the same accounting; the matching-only numbers of
Table~\ref{tab:vitdet-sweep} are not comparable to these and are therefore
kept separate.}
\label{tab:vitdet-sota}
\end{table}

\paragraph{When does content-based reuse help?}
The ViViT and ViTDet results are opposite, and the contrast is the most
informative finding of this evaluation. On classification, the prediction
is read from a class token that aggregates the sequence; replacing a
redundant patch token with a nearby prototype changes what the class token
attends to only marginally, so accuracy is flat while 75\% of the
recompute set is dropped (Table~\ref{tab:eventful-vivit}). On detection,
the prediction is read from the spatial tokens themselves; every dropped
or substituted token is a location the detector can no longer describe
exactly, so mAP@50 falls in proportion to the reuse rate
(Table~\ref{tab:vitdet-sweep}). Two secondary observations support this
reading. First, prototype granularity is inert on classification but
matters on detection, where fewer merge passes (finer prototypes, hence
smaller substitution error per token) are consistently better. Second,
the standalone results follow the same rule: the strongest standalone
numbers come from classification backbones with class-token or pooled
readouts (Tables~\ref{tab:standalone-vivit} and~\ref{tab:scaling}).
TempoMem is therefore best positioned as an accelerator for
classification-style, global-readout video inference, and as a research
direction, rather than a finished method, for dense spatial prediction.

\subsection{Ablation Experiment}
\label{sec:eval-ablations}

\paragraph{Effect of the reuse fraction $\gamma$.}
Covered by Tables~\ref{tab:eventful-vivit} and~\ref{tab:vitdet-sweep}:
a flat step on classification, a monotonic trade on detection.

\paragraph{Effect of prototype granularity ($T$ merge passes).}
On ViViT, $T \in \{2,4,8\}$ is fully inert: identical accuracy and FLOPs
at every $\gamma$. On ViTDet, accuracy degrades with $T$ at every $\gamma$
(e.g., 77.8\% vs.\ 74.7\% mAP@50 at $\gamma = 0.5$), so detection prefers
the finest prototypes the memory budget allows. This is direct evidence
that classification tolerates coarse summaries while detection consumes
the per-token detail that merging removes.

\paragraph{Does content matching help without dropping tokens?}
We also tested a motion-compensated variant that matches stale tokens by
content but never drops them, hypothesizing robustness to large frame
strides. It is refuted: across strides $\{1, 4, 8, 16\}$ on ImageNet VID,
Eventful and the variant degrade identically (differences $\le 0.6$
mAP@50 points, within noise; e.g., 62.9\% vs.\ 62.3\% at stride 8). The
value of the prototype cache lies in \emph{removing} compute, not in
improving the reference features, which motivated the recompute-set filter
as the surviving design.

\paragraph{Warm-up length $W$ and refresh period $P$.}
\textbf{[TBD]}: sweeps over $W$ and $P$ (including scene-cut stress tests)
are planned; the Methodology's failure-mode analysis predicts $P$ to be
the dominant accuracy knob after $r$.

\paragraph{Runtime evaluation: latency and memory.}
FLOPs do not automatically become speed, and we report this honestly.
Table~\ref{tab:vitdet-sota} includes measured wall-clock latency and peak
memory on ImageNet VID. The gated baselines translate their FLOPs savings
into real speedups (56 to 58\,ms vs.\ 69.6\,ms dense). Standalone TempoMem
does not yet: matching and cache management add overhead that outweighs
the saved projections in the current implementation (101.8\,ms), although
its memory footprint (929\,MB) stays close to dense and far below
Eventful's buffers (2118\,MB), reflecting the merged cache of
Proposition~\ref{prop:cachesize}. The gap is implementation-level, not
algorithmic: the matching product and gather/scatter steps are unfused
kernels. Profiling and fused implementations are future work; latency for
the Eventful + TempoMem filter is \textbf{[TBD]}.

\subsection{Performance Across Backbone Scales}
\label{sec:eval-scaling}

Table~\ref{tab:scaling} applies standalone TempoMem to VideoMAE encoders
from ViT-S to ViT-H on Kinetics-400. The relative saving grows with model
size, from 49\% on ViT-S to 86\% on ViT-H, at an accuracy cost between 0.8
and 1.2 points. The reason is
structural: the matching overhead grows as $N_b M C$ while the skipped
projections grow as $r N_b C^2$, so the break-even margin of
Eq.~\ref{eq:breakeven} widens with $C$. Practically, ViT-L with TempoMem
runs at the FLOPs of a dense ViT-S while keeping a 1.8-point Top-1
advantage over it. On UCF101, the same configuration reduces VideoMAE-B
from 156.3 to 86.0 GFLOPs ($-$45\%) with unchanged accuracy (93.4\% vs.\
93.5\%). The results indicate that TempoMem improves computational
efficiency more the larger the backbone, and exhibits favorable
accuracy--computation tradeoffs across all tested scales.

\begin{table}[t]
\centering
\small
\begin{tabular}{llcc}
\toprule
Backbone & Method & GFLOPs & Top-1 (\%) \\
\midrule
ViT-S & dense & 269.2 & 90.5 \\
      & + TempoMem & 137.8 ($-$49\%) & 89.7 \\
\midrule
ViT-B & dense & 502.9 & 91.0 \\
      & + TempoMem & 96.4 ($-$81\%) & 89.8 \\
\midrule
ViT-L & dense & 1459.3 & 93.3 \\
      & + TempoMem & 229.3 ($-$84\%) & 92.3 \\
\midrule
ViT-H & dense & 2895.0 & 94.0 \\
      & + TempoMem & 412.0 ($-$86\%) & 93.1 \\
\bottomrule
\end{tabular}
\caption{Backbone scaling of standalone TempoMem (VideoMAE encoders,
Kinetics-400, per-clip GFLOPs). The relative saving grows with model
capacity, consistent with the break-even analysis: larger $C$ makes
matching cheaper relative to the projections it replaces.}
\label{tab:scaling}
\end{table}

\subsection{Summary of Findings}
\label{sec:eval-summary}

(1)~Standalone TempoMem preserves Top-1 exactly at a 27\% FLOPs reduction
on ViViT and scales with backbone size, reaching 86\% reduction on ViT-H
at a 0.9-point cost; a ViT-L with TempoMem costs less than a dense ViT-S.
(2)~As a filter on Eventful, TempoMem drops up to 75\% of the recompute
set on Kinetics-400 for a fixed 2-point Top-1 cost, cutting matching-frame
compute $3.8\times$ below Eventful itself. (3)~On detection, the same
filter is dominated by Eventful: dense spatial readouts consume exactly
the per-token information that prototype substitution removes, an honest
negative result that scopes the method to global-readout tasks.
(4)~FLOPs savings do not yet translate into latency; memory stays near
dense. Whole-stream measurements for the filter and the $W$/$P$ sweeps
are marked \textbf{[TBD]} throughout.

% ------------------------------------------------------------------
% Additional bib entries used by this section:
% @inproceedings{vivit,   title={ViViT: A Video Vision Transformer},
%   author={Arnab, Anurag and Dehghani, Mostafa and Heigold, Georg and
%           Sun, Chen and Lu{\v{c}}i{\'c}, Mario and Schmid, Cordelia},
%   booktitle={ICCV}, year={2021}}
% @inproceedings{vitdet,  title={Exploring Plain Vision Transformer
%   Backbones for Object Detection},
%   author={Li, Yanghao and Mao, Hanzi and Girshick, Ross and He, Kaiming},
%   booktitle={ECCV}, year={2022}}
% @inproceedings{videomae, title={VideoMAE: Masked Autoencoders are
%   Data-Efficient Learners for Self-Supervised Video Pre-Training},
%   author={Tong, Zhan and Song, Yibing and Wang, Jue and Wang, Limin},
%   booktitle={NeurIPS}, year={2022}}
% @article{k400, title={The Kinetics Human Action Video Dataset},
%   author={Kay, Will and others}, journal={arXiv:1705.06950}, year={2017}}
% @article{ucf101, title={UCF101: A Dataset of 101 Human Actions Classes
%   From Videos in The Wild}, author={Soomro, Khurram and Zamir, Amir
%   Roshan and Shah, Mubarak}, journal={arXiv:1212.0402}, year={2012}}
% @article{imagenet, title={ImageNet Large Scale Visual Recognition
%   Challenge}, author={Russakovsky, Olga and others},
%   journal={IJCV}, year={2015}}
% @inproceedings{maskvd, title={MaskVD: Region Masking for Efficient Video
%   Object Detection}, author={Sarkar, Sreetama and others},
%   year={2024}, booktitle={arXiv:2407.12067}}
% @article{stgt, title={Spatio-Temporal Gated Transformers for Efficient
%   Video Processing}, author={Habibian, Amir and others}, year={2021},
%   journal={NeurIPS Workshops}}
% ------------------------------------------------------------------
